% !PRINTED 655
\weberpage{655}
%<<CONT
so ist $x_1c_1 + x_2c_2 + \cdots + x_nc_n$ durch $\pi$, und weil es eine
Funktion in $Z$ ist, durch $(z - c)$ teilbar. Wir setzen:
\[
x_1c_1 + x_2c_2 + \cdots + x_nc_n = (z - c)y
\]
und erhalten aus (8)
\weq{9}{
\begin{aligned}
\omega &= x_1(\omega_1 - c_1) + x_2(\omega_2 - c_2) + \cdots + x_n(\omega_n - c_n) + (z - c)y\\
&= (\tau - \tau_0) + (z - c)y.
\end{aligned}
}

Hätte nun $(\tau - \tau_0)$ und $(z - c)$ den gemeinschaftlichen Teiler
$\pi\delta$, so könnte man nach \S\ 182, 11. eine durch $\pi$ teilbare, zu $\delta$
teilerfremde Funktion $\omega$ in $\Omega$ finden. Dem widerspricht aber die
Gleichung (9), und damit ist unser Satz bewiesen.

Es sei
\weq{10}{
\Phi(t) = B_0t^m + B_1t^{m-1} + \cdots + B_{m-1}t + B_m
}
eine Funktion $m$ten Grades von $t$, und $B_0, B_1, \ldots, B_m$ seien ganze
holomorphe Funktionale in $\overline{Z}$ mit den Funktionalvariablen $t_1,
t_2, \ldots, t_n$, die der Kongruenzbedingung
\weq{11}{
\Phi(\tau) \equiv 0 \quad (\mathrm{mod}\ z - c)
}
genügt [wie z.\,B. die Funktion $F(t) = N(t - \tau)$]. Dann ist
\[
\Phi(\tau_0) \equiv 0 \quad (\mathrm{mod}\ \pi)\,,
\]
und wenn wir $\Phi(t)$ durch $\tau - \tau_0$ dividieren, so ergibt sich
% [sic] printed $\tau - \tau_0$ (Greek tau) as the divisor of $\Phi(t)$, but division by $t - \tau_0$ is meant, as the factor $(t - \tau_0)$ in (12) shows
\weq{12}{
\Phi(t) \equiv (t - \tau_0)\Phi_1(t) \quad (\mathrm{mod}\ \pi)\,,
}
worin $\Phi_1(t)$ ein holomorphes ganzes Funktional in $\overline{Z}$ ist (weil
$\tau_0$ als \gesperrt{konstantes} Funktional in $\overline{Z}$ enthalten ist). Da nun in (12)
außer den Funktionalvariablen nur rationale Funktionen von $z$
vorkommen, so muß diese Kongruenz auch für den Modul $z - c$
gültig sein:
\weq{13}{
\Phi(t) \equiv (t - \tau_0)\Phi_1(t) \quad (\mathrm{mod}\ z - c).
}

Ist $\pi'$ ein zweites Primfunktional, das auch mit $\pi$ identisch
sein kann, und $z - c$ durch $\pi\pi'$ teilbar, und $\tau \equiv \tau_0' \ (\mathrm{mod}\ \pi')$,
so ist $(\tau - \tau_0)\!: \pi$ relativ prim zu $\pi'$ (nach 3.) und wegen (11)
ist $\Phi_1(\tau_0') \equiv 0 \ (\mathrm{mod}\ \pi')$. Daraus schließt man
\[
\Phi_1(t) \equiv (t - \tau_0')\Phi_2(t)
\]
zunächst für den Modul $\pi'$, dann aber auch, da die Funktionale
in $\overline{Z}$ liegen, auch für den Modul $(z - c)$. Wir haben also:
\weq{14}{
\Phi(t) \equiv (t - \tau_0)(t - \tau_0')\Phi_2(t) \quad (\mathrm{mod}\ z - c).
}

Daraus schließen wir auf folgenden wichtigen Satz:

4. \gesperrt{Es sei nach \S\ 183, (17)}
\weq{15}{
z - c = \pi_1\pi_2\ldots\pi_n
}
%CONT>>
